\documentclass[11pt,a4paper]{amsart}

\usepackage[utf8]{inputenc}
\usepackage[T1]{fontenc}
\usepackage[ngerman]{babel}
\usepackage{amsmath,amssymb,amsthm}
\usepackage{array,booktabs,longtable}
\usepackage{hyperref}

\newtheorem{theorem}{Satz}[section]
\newtheorem{proposition}[theorem]{Proposition}
\newtheorem{lemma}[theorem]{Lemma}
\newtheorem{corollary}[theorem]{Korollar}
\newtheorem{remark}[theorem]{Bemerkung}
\newtheorem{openproblem}[theorem]{Offenes Problem}

\newcommand{\KX}{\mathcal K_X}
\newcommand{\HG}{\mathscr H_{\log}}
\newcommand{\Ran}{\operatorname{Ran}}
\newcommand{\dist}{\operatorname{dist}}

\title[P11 --- Global Coupling]{P11 --- Global Coupling and the Object-X Candidate Geometry}
\author{Objekt-X-Programm}
\date{2026-08-12}

\begin{document}

\begin{abstract}
P11 develops a finite-horizon Hilbert-space geometry for a source-first candidate of
Object~X and isolates, without overclaim, the remaining passage to a global terminal
geometry.  The construction combines a source-windowed archimedean Gamma form with a
finite-horizon Feshbach--Schur contribution and equips the resulting graph spaces with
positive terminal metric operators.  Their normalized transition maps are exact
isometries and satisfy finite-level pullback and cocycle identities.  Relative future
metrics admit an exact compression identity, and a second normalization yields an
isometric modulus map.  On the odd sector, the future Schur term has a sharp fixed-vector
asymptotic of size $e^T/T^{2m+2}$, which forces superpolynomial growth of the associated
fixed-base-terminal condition number.  This makes one natural Jensen-product route to
terminal convergence extremely restrictive and motivates a separate second-moment and
range-leakage diagnostic.  That diagnostic is developed only as a sufficient comparison
mechanism: its failure would not imply failure of strong terminal transport.

The central unresolved source-level problem remains the strong odd terminal limit
\[
 W_{R,S,-}^{[T]}\xrightarrow[T\to\infty]{\rm strong}W_{R,S,-}^{[\infty]}.
\]
P11 does not prove this limit, does not prove its failure, does not construct a final
global Object~X, and contains no proof of the Riemann hypothesis.  The paper therefore
separates three levels throughout: proved finite-horizon structure, route-specific
diagnostics, and open global obligations.
\end{abstract}

\maketitle

\section{Scope, status semantics and non-claims}\label{sec:scope}

The purpose of P11 is to organize the positive finite-horizon geometry that survives the
preceding no-go analysis and to state precisely what is still missing for a global
Object-X construction.  The finite-window theory developed below is mathematically
nontrivial, but it must not be identified with the global limit problem.

\begin{remark}[Three logically distinct levels]\label{rem:three-levels}
We distinguish throughout:
\begin{enumerate}
  \item \emph{finite-horizon structure}: exact Hilbert-space, Feshbach, metric and
        comparison identities at fixed windows;
  \item \emph{terminal transport}: asymptotic compatibility of the normalized maps as
        the terminal horizon tends to infinity;
  \item \emph{P11-wide global geometry}: a canonical nonorthogonal global Gram/mediator
        realization, including the eventual source/adelic and Fredholm/Schatten layers.
\end{enumerate}
A statement at one level is never silently promoted to the next.
\end{remark}

\begin{remark}[Non-claims]\label{rem:nonclaims}
The present draft does not assert any of the following:
\[
W_{R,S,-}^{[T]}\to W_{R,S,-}^{[\infty]}\quad\text{strongly},
\]
\[
\Theta_-\to0,
\qquad
\chi_-\|\Theta_-\|\to0,
\]
any polynomial lower witness for the actual P11 family, any global Fredholm determinant,
a final global Hilbert space $\KX$, or the Riemann hypothesis.
\end{remark}

\begin{remark}[Draft provenance]\label{rem:provenance}
This first paper draft is distilled from the canonical P11 master ledger.  Historical
audit nodes remain proof provenance, but superseded formulations are not treated as
independent inputs.  In particular, the full-rest transfer is used only in its repaired
martingale form, the old primitive form-domination shortcut is not assumed, and the
O3 diagnostic branch is kept separate from the original strong transport target.
\end{remark}

\section{Fixed-window source geometry}\label{sec:source}

Fix $R>0$ and write
\[
  \mathscr H_R:=L^2(-R,R).
\]
For $0<R<S$, let
\[
  J_{R,S}:\mathscr H_R\longrightarrow\mathscr H_S
\]
denote zero extension.  If a common terminal horizon $T>S$ is fixed, we also write
$E_R=J_{R,T}$ and $E_S=J_{S,T}$.

The archimedean part is encoded by the positive Fourier weight
\[
 m_\Gamma(\xi):=1+g_\infty(\xi),
\]
where the committed source-windowed Gamma construction satisfies
\[
  m_\Gamma(\xi)\asymp \log(2+|\xi|).
\]
The corresponding closed form on $\mathscr H_R$ is
\[
 \mathfrak c_{\Gamma,R}[f,g]
 =\frac1{2\pi}\int_{\mathbb R}
   m_\Gamma(\xi)
   \widehat{E_Rf}(\xi)
   \overline{\widehat{E_Rg}(\xi)}\,d\xi.
\]
We denote its positive selfadjoint operator by $C_{\Gamma,R}$.

\begin{proposition}[Gamma graph space]\label{prop:gamma-graph}
For every fixed $R>0$, the form $\mathfrak c_{\Gamma,R}$ is densely defined, closed and
positive.  Its form domain, equipped with the graph norm, defines the native source
space
\[
 \KX_R:=\bigl(\mathcal D(\mathfrak c_{\Gamma,R}),\|\cdot\|_{X,R}\bigr).
\]
Moreover $C_{\Gamma,R}\ge I$.  Since $m_\Gamma(\xi)\to\infty$, the fixed-window graph
embedding into $L^2(-R,R)$ is compact, hence the finite-window Gamma realization has
compact resolvent.
\end{proposition}

\begin{remark}[Logarithmic, not Sobolev, regularity]\label{rem:log-reg}
The native form norm corresponds to the logarithmic threshold
\[
 \mathcal D(\mathfrak c_{\Gamma,R})
 \simeq
 \left\{f:\int \log(2+|\xi|)\,|\widehat{E_Rf}(\xi)|^2\,d\xi<\infty\right\}.
\]
It does not force membership of the entire form domain in any fixed positive Sobolev
space $H^s$, $s>0$.
\end{remark}

\section{Finite-horizon Feshbach geometry}\label{sec:feshbach}

At terminal horizon $T>R$, the source couples to a hub/rest system.  The exact
finite-horizon Schur contribution has the operator form
\[
 \Sigma_T
 =H_T\,(I+R_T^*R_T)^{-1}H_T^*\ge0.
\]
The source-restricted terminal form is therefore
\[
 q_T^X(E_Rf,E_Rg)
 =\mathfrak c_{\Gamma,R}[f,g]
 +\langle \Sigma_R^{[T]}f,g\rangle_{L^2(-R,R)},
\]
where
\[
 \Sigma_R^{[T]}:=E_R^*\Sigma_T E_R
\]
is bounded, positive and selfadjoint.

\begin{proposition}[Bounded Schur perturbation]\label{prop:schur-bounded}
For every fixed $R<T$, the operator
\[
 F_R^{[T]}:=C_{\Gamma,R}+\Sigma_R^{[T]}
\]
is positive selfadjoint with
\[
 \mathcal D(F_R^{[T]})=\mathcal D(C_{\Gamma,R}).
\]
In particular the finite-horizon Schur term does not alter the operator domain; it is a
bounded positive perturbation of the Dirichlet Gamma realization.
\end{proposition}

\begin{proof}
The operator $\Sigma_R^{[T]}$ is bounded and selfadjoint because
$(I+R_T^*R_T)^{-1}$ is a positive contraction and $H_T$ is bounded at fixed $T$.
The standard bounded selfadjoint perturbation theorem gives the domain identity and
selfadjointness of $F_R^{[T]}$.
\end{proof}

\subsection{Full-rest martingale representation}

The rest Gram is not a direct orthogonal sum over prime powers $k$ inside a fixed prime
sector.  The correct positive decomposition is instead indexed by prime/martingale-depth
pairs $(p,a)$.  For
\[
 \Phi_{p,a,T}[f](u)
 :=\sum_{k\ge a+1}p^{-3k/4}K_{k\log p}f(u)
\]
and the corresponding source cells $\Omega_{p,a,T}$, the residual bilinear form is
\[
 \mathfrak R_T(f,g)
 =\sum_p(\log p)(p-1)\sum_{a\ge0}p^a
   \int_{\Omega_{p,a,T}}
   \Phi_{p,a,T}[f](u)\overline{\Phi_{p,a,T}[g](u)}\,du.
\]
Thus there is a canonical analysis operator $\widetilde R_T$ satisfying
\[
 \boxed{\widetilde R_T^*\widetilde R_T=R_T^*R_T.}
\]

\begin{remark}[Primitive form-domination firewall]\label{rem:primitive-fw}
No use is made of the unproved global form inequality
\[
 R_T^*R_T\stackrel?\ge (R_T^{(1)})^*R_T^{(1)}.
\]
Prime powers inside one prime sector are nonorthogonal, so that inequality is not an
automatic consequence of channel decomposition.  The required future-certificate
transfer is instead obtained through the full martingale analysis operator
$\widetilde R_T$ and the exponentially small higher-prime-power tail in the future
$a=0$ channel.
\end{remark}

\section{Terminal metric operators and normalized transports}\label{sec:terminal}

The fixed-window source spaces carry their native graph inner products.  For $R<T$, let
\[
 J_{R,T}:\KX_R\longrightarrow\KX_T
\]
be the terminal embedding.  Define the terminal metric operator
\[
 \boxed{G_{R,T}:=J_{R,T}^*J_{R,T}\in\mathcal B(\KX_R).}
\]
It is positive and boundedly invertible.

\begin{proposition}[Pullback identity]\label{prop:pullback}
For $0<R<S<T$,
\[
 \boxed{J_{R,S}^*G_{S,T}J_{R,S}=G_{R,T}.}
\]
\end{proposition}

\begin{proof}
Using $J_{S,T}J_{R,S}=J_{R,T}$,
\[
 J_{R,S}^*G_{S,T}J_{R,S}
 =J_{R,S}^*J_{S,T}^*J_{S,T}J_{R,S}
 =J_{R,T}^*J_{R,T}=G_{R,T}.
\]
\end{proof}

The normalized terminal transport is
\[
 \boxed{
 W_{R,S}^{[T]}
 :=G_{S,T}^{1/2}J_{R,S}G_{R,T}^{-1/2}.
 }
\]

\begin{theorem}[Finite-terminal isometry]\label{thm:terminal-isometry}
For every $0<R<S<T$,
\[
 \boxed{(W_{R,S}^{[T]})^*W_{R,S}^{[T]}=I.}
\]
Hence $W_{R,S}^{[T]}$ is an isometry.  The family satisfies the compatible
finite-terminal cocycle identities whenever the windows are nested.
\end{theorem}

\begin{proof}
By Proposition~\ref{prop:pullback},
\[
\begin{aligned}
 (W_{R,S}^{[T]})^*W_{R,S}^{[T]}
 &=G_{R,T}^{-1/2}J_{R,S}^*G_{S,T}J_{R,S}G_{R,T}^{-1/2}\\
 &=I.
\end{aligned}
\]
The cocycle follows from the same pullback identity and functoriality of the nested
zero-extension maps.
\end{proof}

\begin{remark}[Finite isometry is not terminal convergence]\label{rem:not-limit}
The exact identity in Theorem~\ref{thm:terminal-isometry} is a fixed-$T$ statement.  It
does not imply that the family $W_{R,S}^{[T]}$ is Cauchy as $T\to\infty$.
\end{remark}

\section{Relative future metrics}\label{sec:relative}

Fix
\[
 0<R<S<T<U.
\]
On the $R$- and $S$-source spaces define the relative future metrics
\[
 A_R^{T,U}:=G_{R,T}^{-1/2}G_{R,U}G_{R,T}^{-1/2},
\]
\[
 A_S^{T,U}:=G_{S,T}^{-1/2}G_{S,U}G_{S,T}^{-1/2}.
\]
Write $W=W_{R,S}^{[T]}$.

\begin{theorem}[Exact relative compression]\label{thm:relative-compression}
The future metrics satisfy
\[
 \boxed{W^*A_S^{T,U}W=A_R^{T,U}.}
\]
\end{theorem}

\begin{proof}
Insert the definitions and use twice the pullback identity
$J_{R,S}^*G_{S,V}J_{R,S}=G_{R,V}$, first for $V=T$ and then for $V=U$.
\end{proof}

Define the second normalized map
\[
 Q:=\bigl(A_S^{T,U}\bigr)^{1/2}
 W\bigl(A_R^{T,U}\bigr)^{-1/2}.
\]

\begin{corollary}[Modulus isometry]\label{cor:modulus-isometry}
\[
 \boxed{Q^*Q=I.}
\]
\end{corollary}

\begin{proof}
This is immediate from Theorem~\ref{thm:relative-compression}.
\end{proof}

\begin{remark}[What the first-power identity does not give]\label{rem:first-power}
The compression
\[
 W^*A_SW=A_R
\]
does not by itself imply invariance of $\Ran W$ under $A_S$, nor does it determine the
compression of $A_S^{1/2}$.  This is the origin of the Jensen/range defect studied below.
\end{remark}

\section{Odd future asymptotics and conditioning}\label{sec:odd}

The odd sector detects a sharp future growth mechanism.  Let
\[
 f_-\in C_c^\infty((-R,R)),\qquad f_-(-u)=-f_-(u),
\]
and let
\[
 m:=\min\{j\ge0:\beta_R^{(j)}(f_-)\ne0\}
\]
be the first nonvanishing integral jet.  Here the functionals
$\beta_R^{(j)}$ are the canonical integral jets of the source kernel expansion; they are
not local boundary derivatives.

Set
\[
 c_m:=\binom{2m}{m}4^{-m}.
\]

\begin{theorem}[Sharp fixed-vector odd future asymptotic]\label{thm:odd-asymptotic}
For fixed $R$ and fixed nonzero smooth odd $f_-$ as above,
\[
 \boxed{
 \sigma_T(J_{R,T}f_-)
 =c_m^2|\beta_R^{(m)}(f_-)|^2
  \frac{e^T}{T^{2m+2}}\bigl(1+o(1)\bigr).
 }
\]
The proof uses the exact constant-mode denominator
\[
 \langle\mathbf1_T,(I+R_T^*R_T)\mathbf1_T\rangle=2T+O(1),
\]
a signed mean-zero future-edge certificate, prime-cell quadrature, and the full-rest
martingale lift from Section~\ref{sec:feshbach}.
\end{theorem}

\begin{remark}[Proof status in Draft 0.1]\label{rem:odd-proof}
The theorem is an audited result, but this first paper draft records only the proof
architecture.  The final P11 proof will expand the constant-mode estimate, the signed
certificate cost, the prime quadrature and the full-rest squeeze into self-contained
lemmas rather than cite the historical audit chain.
\end{remark}

Fix now a base terminal horizon $T_0>R$.  The relative odd metric on the fixed source
space has Rayleigh quotient
\[
 \rho_{T_0,U}(f)
 =\frac{\langle G_{R,U}f,f\rangle}
        {\langle G_{R,T_0}f,f\rangle}.
\]
Theorem~\ref{thm:odd-asymptotic} yields, for the fixed vector $f_-$,
\[
 \rho_{T_0,U}(f_-)
 \sim C_{T_0,f_-}\frac{e^U}{U^{2m+2}}.
\]
Since the order $m$ can be chosen arbitrarily large within the odd test core, the
condition number of the odd relative metric obeys the following strong divergence.

\begin{corollary}[Superpolynomial odd conditioning]\label{cor:conditioning}
For every $N>0$,
\[
 \boxed{
 U^{-N}\kappa\bigl(A_{T_0,U}^{R,-}\bigr)\longrightarrow\infty.
 }
\]
Consequently, for
\[
 \chi_{T_0,U}^{R,-}
 :=\kappa\bigl(A_{T_0,U}^{R,-}\bigr)^{1/4},
\]
one also has
\[
 \boxed{
 \forall N>0:\qquad
 U^{-N}\chi_{T_0,U}^{R,-}\longrightarrow\infty.
 }
\]
\end{corollary}

\section{Jensen geometry as a sufficient diagnostic route}\label{sec:jensen}

Fix $0<R<S<T_0<U$ and restrict to odd sectors.  Write
\[
 A_R:=A_{T_0,U}^{R,-},\qquad
 A_S:=A_{T_0,U}^{S,-},\qquad
 W:=W_{R,S,-}^{[T_0]}.
\]
Then
\[
 W^*A_SW=A_R,
\qquad W^*W=I.
\]
Define
\[
 D:=W^*A_S^{1/2}W,
\qquad
 \mathscr J:=A_R^{1/2}-D\ge0,
\]
and the symmetrized Jensen defect
\[
 \boxed{
 \Theta:=A_R^{-1/4}\mathscr J A_R^{-1/4},
 \qquad 0\le\Theta\le I.
 }
\]
Further set
\[
 \chi:=\|A_R^{1/4}\|\,\|A_R^{-1/4}\|
 =\kappa(A_R)^{1/4}.
\]

\begin{proposition}[Relative Jensen Rayleigh formula]\label{prop:jensen-rayleigh}
\[
 \boxed{
 \|\Theta\|
 =\sup_{y\ne0}
 \left[
 1-
 \frac{\langle A_S^{1/2}Wy,Wy\rangle}
      {\langle A_R^{1/2}y,y\rangle}
 \right].
 }
\]
\end{proposition}

Let $P:=WW^*$, which is the orthogonal projection onto the closed range of the isometry
$W$, and define
\[
 Q:=A_S^{1/2}WA_R^{-1/2},
\quad
 \mathscr N:=(I-P)Q,
\quad
 \mathscr K:=I-W^*Q.
\]
The exact defect balance is
\[
 \boxed{
 \mathscr K+\mathscr K^*
 =\mathscr K^*\mathscr K+\mathscr N^*\mathscr N.
 }
\]
Moreover
\[
 \mathscr K=A_R^{1/4}\Theta A_R^{-1/4},
\]
so that
\[
 \boxed{
 \frac12\|\mathscr N\|^2
 \le\|\mathscr K\|
 \le\chi\|\Theta\|.
 }
\]

\begin{corollary}[Beyond-all-orders necessity for this route]\label{cor:beyond}
If the sufficient product mechanism
\[
 \chi_{T_0,U}^{R,-}\|\Theta_{T_0,U}^-\|\to0
\]
holds, then for every $N>0$,
\[
 \boxed{U^N\|\Theta_{T_0,U}^-\|\to0.}
\]
Thus this particular route can succeed only if the odd Jensen defect decays faster than
every inverse power of $U$.
\end{corollary}

\begin{remark}[Critical logical firewall]\label{rem:product-fw}
Corollary~\ref{cor:beyond} concerns a \emph{sufficient comparison mechanism}.  The
present theory does not prove
\[
 \text{strong terminal transport}
 \Longrightarrow
 \chi\|\Theta\|\to0.
\]
Therefore a future obstruction to the product mechanism would not be a proof of
nonconvergence of $W_{R,S,-}^{[T]}$.
\end{remark}

\section{Second-moment and complement diagnostics}\label{sec:diagnostic}

The missing information in the Jensen problem is genuine range geometry.  Define
\[
 \mathscr B:=(I-P)A_SW.
\]
Then
\[
 \mathscr B=0
 \quad\Longleftrightarrow\quad
 A_S\Ran W\subseteq\Ran W.
\]
Since $A_S$ is selfadjoint, this is the reducing-subspace case.

\begin{proposition}[Second-moment compression variance]\label{prop:second-moment}
\[
 \boxed{
 \mathscr B^*\mathscr B
 =W^*A_S^2W-A_R^2
 =:\Delta_2\ge0.
 }
\]
Moreover
\[
 \boxed{
 \|\Theta\|
 \ge
 \frac{\|\Delta_2\|}{8\|A_R\|\|A_S\|}.
 }
\]
\end{proposition}

Define the dimensionless second-moment witness
\[
 \nu_2
 :=\frac{\|\Delta_2\|}{\|A_R\|\|A_S\|}.
\]
A polynomial lower bound for the actual family would therefore contradict the
beyond-all-orders requirement of Corollary~\ref{cor:beyond}.  Such a lower bound is
\emph{not} proved in the present paper.

The exact future cross-Gram representation uses the fixed $T_0$ complement
\[
 \mathcal C^-_{S,T_0}(R)
 :=\ker(J_{R,S}^*G_{S,T_0})\cap\KX_S^-.
\]
For $f\in\KX_R^-$ and $g\in\mathcal C^-_{S,T_0}(R)$,
\[
 \langle\mathscr By,z\rangle
 =\langle(G_{S,U}-G_{S,T_0})J_{R,S}f,g\rangle
\]
after the natural metric normalizations.  This identifies the missing second-moment
information with a scalar future cross-Gram between the old range and its fixed-terminal
Gram complement.

\subsection{The explicit rough complement and the logarithmic gate}

Let $h\in C_c^\infty((-S,S))$ be a nonzero smooth odd annulus vector supported away from
$[-R,R]$.  With
\[
 \Pi^{\rm raw}
 :=J_{R,S}G_{R,T_0}^{-1}J_{R,S}^*G_{S,T_0},
\]
set
\[
 \boxed{g_h:=(I-\Pi^{\rm raw})h.}
\]
Then
\[
 g_h\in\mathcal C^-_{S,T_0}(R),
 \qquad g_h\ne0,
\]
and $g_h$ is compactly supported.  Its first nonzero integral jet has finite order
$m_h$.

For
\[
 \HG^\alpha(\mathbb R)
 :=\left\{F\in L^2:
 \int [\log(2+|\xi|)]^{2\alpha}|\widehat F(\xi)|^2d\xi<\infty
 \right\},
\]
one has the translation estimate
\[
 F\in\HG^\alpha
 \quad\Longrightarrow\quad
 \|\tau_hF-F\|_2
 =o\bigl((\log(1/h))^{-\alpha}\bigr).
\]
The robust prime-cell quadrature therefore closes for a compact odd vector with first
jet $m$ once
\[
 E_Sf\in\HG^{m+3/2}.
\]
The native Gamma form supplies only the baseline exponent $1/2$; the required additional
logarithmic order is $m+1$.

\begin{openproblem}[Logarithmic complement regularity]\label{open:log-complement}
Does the explicit complement vector $g_h$ satisfy
\[
 E_Sg_h\in\HG^{m_h+3/2},
\]
or at least the corresponding direct translation-modulus condition on the prime-cell
scale?  A positive answer would produce a polynomial second-moment witness and hence
rule out the current Jensen-product sufficiency route.  It would still not rule out
strong terminal transport itself.
\end{openproblem}

\subsection{Dirichlet--Riesz reduction}

Write
\[
 u_h:=G_{R,T_0}^{-1}J_{R,S}^*G_{S,T_0}h,
\qquad
 g_h=h-J_{R,S}u_h.
\]
Then $J_{R,T_0}u_h$ is the orthogonal projection of $J_{S,T_0}h$ onto
$\Ran J_{R,T_0}$ in the $T_0$ Object-X metric.  Equivalently,
\[
 q_{T_0}^X(E_Ru_h,E_Rv)
 =q_{T_0}^X(E_Sh,E_Rv)
 \qquad(v\in\KX_R).
\]
The forcing decomposes into a Gamma part and a Schur part,
\[
 r_h=r_{\Gamma,h}+r_{\sigma,h}.
\]
For smooth interior $h$, the Gamma part has arbitrary finite logarithmic regularity,
whereas the committed theory gives only $L^2$ control for the Schur forcing
$r_{\sigma,h}$.  Since the Schur term is bounded,
\[
 \boxed{
 u_h
 =(C_{\Gamma,R}+\Sigma_R^{[T_0]})^{-1}r_h
 \in\mathcal D(C_{\Gamma,R}).
 }
\]
No theorem in the present draft upgrades this operator-domain gain to the threshold in
Open Problem~\ref{open:log-complement}.

\begin{remark}[Dirichlet Fourier firewall]\label{rem:dirichlet-fw}
The metric inverse $G_{R,T_0}^{-1}$ is not the global Fourier multiplier
$(1+g_\infty)^{-1}$.  The finite-window source compression is a genuine Dirichlet/Riesz
problem and must not be replaced by a global symbol inversion.
\end{remark}

\section{The open strong terminal transport problem}\label{sec:open-transport}

The finite-terminal maps of Theorem~\ref{thm:terminal-isometry} form the central
source-level transport family.  The original asymptotic target is
\[
 \boxed{
 W_{R,S,-}^{[T]}
 \xrightarrow[T\to\infty]{\rm strong}?
 W_{R,S,-}^{[\infty]}.
 }
\]

\begin{openproblem}[Strong odd terminal transport]\label{open:transport}
For fixed $0<R<S$, is the family $W_{R,S,-}^{[T]}$ strongly Cauchy as the common terminal
horizon $T\to\infty$?  Equivalently, can the source-level cross-terminal Cauchy kernel be
controlled on the full odd graph space?
\end{openproblem}

\begin{remark}[Two research directions]\label{rem:two-directions}
The current theory leaves two distinct directions:
\begin{enumerate}
  \item a \emph{direct transport route}, attacking the cross-terminal Cauchy asymptotic
        itself;
  \item the \emph{O3 diagnostic route}, studying whether the sufficient
        Jensen-product mechanism is compatible with the superpolynomial odd
        conditioning.
\end{enumerate}
The second route is not a semantic replacement for the first.
\end{remark}

Residual observability, signed/clustered finite-band estimates and related quantities
may contribute to a direct route, but no equivalence between those diagnostics and
strong terminal transport is asserted here.

\section{Global obligations beyond the finite-horizon core}\label{sec:global}

Even a solution of Open Problem~\ref{open:transport} would not by itself complete the
full P11 program.  The historical P11 scope also carries genuinely global obligations.

\begin{openproblem}[Global nonorthogonal Gram/mediator closure]\label{open:gram}
Construct a canonical global positive geometry in which the local prime-power channels
and the archimedean contribution arise from one nonorthogonal source/mediator mechanism,
rather than from an external additive juxtaposition.
\end{openproblem}

\begin{openproblem}[Canonical source/adelic realization]\label{open:adelic}
Identify a global source space and canonical realization compatible with the finite
source-window geometry and with the adelic/Weil amplitude port, including the correct
prime-power and archimedean normalization.
\end{openproblem}

\begin{openproblem}[Global Fredholm/Schatten layer]\label{open:fredholm}
Determine whether the finite-horizon Feshbach operators admit a canonical global limit
with the Schatten/Fredholm control required for a genuine determinant or resolvent
formalism.  Finite Feshbach identities alone do not provide such a limit.
\end{openproblem}

\section{Conclusion}\label{sec:conclusion}

The finite-horizon Object-X candidate geometry is substantially more rigid than a mere
collection of local channels.  It provides a closed logarithmic Gamma graph space,
positive Feshbach corrections, exact terminal pullback metrics, normalized transition
isometries, relative future compression, a second modulus isometry, and sharp odd
fixed-vector future asymptotics.  These results are compatible with a source-first
nonorthogonal geometry and survive the earlier no-go analysis.

At the same time, the main asymptotic bottleneck is explicit.  Strong odd terminal
transport remains open.  The superpolynomial odd conditioning shows that one natural
Jensen-product sufficiency route would require beyond-all-orders decay of a square-root
compression defect.  The second-moment/complement analysis reduces a possible
obstruction to a concrete logarithmic regularity problem, but that problem is still
open and, even if resolved negatively for the product route, would not decide the
original strong transport question.

Accordingly, the mathematically correct current status is:
\[
 \boxed{
 \text{finite-horizon structural core: proved;}
 \quad
 \text{strong odd terminal transport: open;}
 \quad
 \text{P11-wide global closure: open.}
 }
\]

\end{document}
